\documentclass[reqno]{amsart}
\usepackage[T1]{fontenc}
\usepackage{amsmath,amssymb,amsthm}
\usepackage{xr-hyper}
\usepackage[hidelinks]{hyperref}
\externaldocument[di-]{20261125T020100--divisibility-at-inert-primes-in-general__workbook}
\externaldocument[nk-]{20261126T020110--the-normalized-kernel-law__workbook}

\newcounter{chap}
% Chapter numbers, generated from 20261125T012000--grh-workbook__guide.org
% by scripts/make-chapters.py.  Do not edit: change the order of the
% index and rerun the script (./preview-site.sh and CI do this).
%
% Usage in a sheet, after \newcounter{chap}:
%   % Chapter numbers, generated from 20261125T012000--grh-workbook__guide.org
% by scripts/make-chapters.py.  Do not edit: change the order of the
% index and rerun the script (./preview-site.sh and CI do this).
%
% Usage in a sheet, after \newcounter{chap}:
%   % Chapter numbers, generated from 20261125T012000--grh-workbook__guide.org
% by scripts/make-chapters.py.  Do not edit: change the order of the
% index and rerun the script (./preview-site.sh and CI do this).
%
% Usage in a sheet, after \newcounter{chap}:
%   \input{chapters}\chapnumber{<sheet ID>}
\makeatletter
\providecommand\chapters@set[2]{\@namedef{chapters@#1}{#2}}
\chapters@set{20261125T001000}{1}% zeros-to-ramp-sums
\chapters@set{20261125T001100}{2}% inert-primes-and-estimate-b
\chapters@set{20261125T005000}{3}% galois-fourier-and-averaging
\chapters@set{20261125T005500}{4}% columns-from-a-sphere
\chapters@set{20261125T000000}{5}% slope-method
\chapters@set{20261126T030100}{6}% a-twist-for-each-prime
\chapters@set{20261125T003000}{7}% the-chain-and-the-square-root
\chapters@set{20261125T010000}{8}% rank-lemma-and-zero-estimate
\chapters@set{20261125T010500}{9}% poincare-series-and-near-orthogonality
\chapters@set{20261124T223000}{10}% elliptic-curves-for-the-proof
\chapters@set{20261125T009500}{11}% fock-model-and-theta-intertwiners
\chapters@set{20261124T223100}{12}% complex-multiplication-and-inert-reduction
\chapters@set{20261125T008000}{13}% one-cm-determinant
\chapters@set{20261124T214500}{14}% supersingular-theta-divisibility
\chapters@set{20261124T214600}{15}% the-isotropic-plane
\chapters@set{20261125T013700}{16}% pairings-on-a-formal-group-of-height-two
\chapters@set{20261125T013500}{17}% hecke-nondegeneracy-on-the-tree
\chapters@set{20261126T030200}{18}% the-bound-is-attained
\chapters@set{20261125T009000}{19}% kronecker-limit-and-chowla-selberg
\chapters@set{20261125T002000}{20}% definite-theta-lifts
\chapters@set{20261125T015500}{21}% siegel-weil-in-the-smallest-case
\chapters@set{20261125T002100}{22}% definite-theta-lifts-and-norms
\chapters@set{20261125T016500}{23}% hecke-operators-and-theta-lifts
\chapters@set{20261125T016700}{24}% divisibility-along-the-chain
\chapters@set{20261125T004000}{25}% denominators-and-sha
\chapters@set{20261125T004100}{26}% tate-twists-and-sha
\chapters@set{20261126T020430}{27}% the-minus-part-by-gauss-sums
\chapters@set{20261125T014100}{28}% congruence-modules
\chapters@set{20261125T014500}{29}% eisenstein-congruences-and-selmer-classes
\chapters@set{20261125T014700}{30}% pseudocharacters-and-the-length-bound
\chapters@set{20261125T011000}{31}% galois-representations-black-box
\chapters@set{20261126T030300}{32}% the-data
\chapters@set{20261126T020210}{33}% bessel-poincare-series-and-their-gram-operators
\chapters@set{20261126T030400}{34}% the-seesaw-identity
\chapters@set{20261126T030500}{35}% divisibility-of-the-composite-lift
\chapters@set{20261126T030600}{36}% theorem-a-short-forms
\chapters@set{20261126T030700}{37}% the-determinant-of-a-chain
\chapters@set{20261126T030800}{38}% the-cost-of-one-lift
\chapters@set{20261126T030900}{39}% denominators-of-projections-to-lifts
\chapters@set{20261126T031000}{40}% twists-with-few-congruences
\chapters@set{20261126T031100}{41}% theorem-c-and-the-main-inequality
\providecommand\chapnumber[1]{%
  \@ifundefined{chapters@#1}%
    {\PackageError{chapters}{Sheet #1 is not listed in the index}%
      {List it in 20261125T012000--grh-workbook__guide.org\MessageBreak
       and run scripts/make-chapters.py.}}%
    {\edef\thechap{\@nameuse{chapters@#1}}}}
\makeatother
\chapnumber{<sheet ID>}
\makeatletter
\providecommand\chapters@set[2]{\@namedef{chapters@#1}{#2}}
\chapters@set{20261125T001000}{1}% zeros-to-ramp-sums
\chapters@set{20261125T001100}{2}% inert-primes-and-estimate-b
\chapters@set{20261125T005000}{3}% galois-fourier-and-averaging
\chapters@set{20261125T005500}{4}% columns-from-a-sphere
\chapters@set{20261125T000000}{5}% slope-method
\chapters@set{20261126T030100}{6}% a-twist-for-each-prime
\chapters@set{20261125T003000}{7}% the-chain-and-the-square-root
\chapters@set{20261125T010000}{8}% rank-lemma-and-zero-estimate
\chapters@set{20261125T010500}{9}% poincare-series-and-near-orthogonality
\chapters@set{20261124T223000}{10}% elliptic-curves-for-the-proof
\chapters@set{20261125T009500}{11}% fock-model-and-theta-intertwiners
\chapters@set{20261124T223100}{12}% complex-multiplication-and-inert-reduction
\chapters@set{20261125T008000}{13}% one-cm-determinant
\chapters@set{20261124T214500}{14}% supersingular-theta-divisibility
\chapters@set{20261124T214600}{15}% the-isotropic-plane
\chapters@set{20261125T013700}{16}% pairings-on-a-formal-group-of-height-two
\chapters@set{20261125T013500}{17}% hecke-nondegeneracy-on-the-tree
\chapters@set{20261126T030200}{18}% the-bound-is-attained
\chapters@set{20261125T009000}{19}% kronecker-limit-and-chowla-selberg
\chapters@set{20261125T002000}{20}% definite-theta-lifts
\chapters@set{20261125T015500}{21}% siegel-weil-in-the-smallest-case
\chapters@set{20261125T002100}{22}% definite-theta-lifts-and-norms
\chapters@set{20261125T016500}{23}% hecke-operators-and-theta-lifts
\chapters@set{20261125T016700}{24}% divisibility-along-the-chain
\chapters@set{20261125T004000}{25}% denominators-and-sha
\chapters@set{20261125T004100}{26}% tate-twists-and-sha
\chapters@set{20261126T020430}{27}% the-minus-part-by-gauss-sums
\chapters@set{20261125T014100}{28}% congruence-modules
\chapters@set{20261125T014500}{29}% eisenstein-congruences-and-selmer-classes
\chapters@set{20261125T014700}{30}% pseudocharacters-and-the-length-bound
\chapters@set{20261125T011000}{31}% galois-representations-black-box
\chapters@set{20261126T030300}{32}% the-data
\chapters@set{20261126T020210}{33}% bessel-poincare-series-and-their-gram-operators
\chapters@set{20261126T030400}{34}% the-seesaw-identity
\chapters@set{20261126T030500}{35}% divisibility-of-the-composite-lift
\chapters@set{20261126T030600}{36}% theorem-a-short-forms
\chapters@set{20261126T030700}{37}% the-determinant-of-a-chain
\chapters@set{20261126T030800}{38}% the-cost-of-one-lift
\chapters@set{20261126T030900}{39}% denominators-of-projections-to-lifts
\chapters@set{20261126T031000}{40}% twists-with-few-congruences
\chapters@set{20261126T031100}{41}% theorem-c-and-the-main-inequality
\providecommand\chapnumber[1]{%
  \@ifundefined{chapters@#1}%
    {\PackageError{chapters}{Sheet #1 is not listed in the index}%
      {List it in 20261125T012000--grh-workbook__guide.org\MessageBreak
       and run scripts/make-chapters.py.}}%
    {\edef\thechap{\@nameuse{chapters@#1}}}}
\makeatother
\chapnumber{<sheet ID>}
\makeatletter
\providecommand\chapters@set[2]{\@namedef{chapters@#1}{#2}}
\chapters@set{20261125T001000}{1}% zeros-to-ramp-sums
\chapters@set{20261125T001100}{2}% inert-primes-and-estimate-b
\chapters@set{20261125T005000}{3}% galois-fourier-and-averaging
\chapters@set{20261125T005500}{4}% columns-from-a-sphere
\chapters@set{20261125T000000}{5}% slope-method
\chapters@set{20261126T030100}{6}% a-twist-for-each-prime
\chapters@set{20261125T003000}{7}% the-chain-and-the-square-root
\chapters@set{20261125T010000}{8}% rank-lemma-and-zero-estimate
\chapters@set{20261125T010500}{9}% poincare-series-and-near-orthogonality
\chapters@set{20261124T223000}{10}% elliptic-curves-for-the-proof
\chapters@set{20261125T009500}{11}% fock-model-and-theta-intertwiners
\chapters@set{20261124T223100}{12}% complex-multiplication-and-inert-reduction
\chapters@set{20261125T008000}{13}% one-cm-determinant
\chapters@set{20261124T214500}{14}% supersingular-theta-divisibility
\chapters@set{20261124T214600}{15}% the-isotropic-plane
\chapters@set{20261125T013700}{16}% pairings-on-a-formal-group-of-height-two
\chapters@set{20261125T013500}{17}% hecke-nondegeneracy-on-the-tree
\chapters@set{20261126T030200}{18}% the-bound-is-attained
\chapters@set{20261125T009000}{19}% kronecker-limit-and-chowla-selberg
\chapters@set{20261125T002000}{20}% definite-theta-lifts
\chapters@set{20261125T015500}{21}% siegel-weil-in-the-smallest-case
\chapters@set{20261125T002100}{22}% definite-theta-lifts-and-norms
\chapters@set{20261125T016500}{23}% hecke-operators-and-theta-lifts
\chapters@set{20261125T016700}{24}% divisibility-along-the-chain
\chapters@set{20261125T004000}{25}% denominators-and-sha
\chapters@set{20261125T004100}{26}% tate-twists-and-sha
\chapters@set{20261126T020430}{27}% the-minus-part-by-gauss-sums
\chapters@set{20261125T014100}{28}% congruence-modules
\chapters@set{20261125T014500}{29}% eisenstein-congruences-and-selmer-classes
\chapters@set{20261125T014700}{30}% pseudocharacters-and-the-length-bound
\chapters@set{20261125T011000}{31}% galois-representations-black-box
\chapters@set{20261126T030300}{32}% the-data
\chapters@set{20261126T020210}{33}% bessel-poincare-series-and-their-gram-operators
\chapters@set{20261126T030400}{34}% the-seesaw-identity
\chapters@set{20261126T030500}{35}% divisibility-of-the-composite-lift
\chapters@set{20261126T030600}{36}% theorem-a-short-forms
\chapters@set{20261126T030700}{37}% the-determinant-of-a-chain
\chapters@set{20261126T030800}{38}% the-cost-of-one-lift
\chapters@set{20261126T030900}{39}% denominators-of-projections-to-lifts
\chapters@set{20261126T031000}{40}% twists-with-few-congruences
\chapters@set{20261126T031100}{41}% theorem-c-and-the-main-inequality
\providecommand\chapnumber[1]{%
  \@ifundefined{chapters@#1}%
    {\PackageError{chapters}{Sheet #1 is not listed in the index}%
      {List it in 20261125T012000--grh-workbook__guide.org\MessageBreak
       and run scripts/make-chapters.py.}}%
    {\edef\thechap{\@nameuse{chapters@#1}}}}
\makeatother

\ifcsname chapters@20261126T020120\endcsname
  \chapnumber{20261126T020120}
\else
  \renewcommand{\thechap}{?}
\fi
\renewcommand{\thesection}{\thechap.\arabic{section}}
\theoremstyle{definition}
\newtheorem{problem}{Problem}[section]
\theoremstyle{plain}
\newtheorem{theorem}{Theorem}[section]
\theoremstyle{remark}
\newtheorem*{leanremark}{In the Lean formalization}
\newcommand{\Z}{\mathbb Z}
\newcommand{\Q}{\mathbb Q}
\newcommand{\Oc}{\mathcal O}
\newcommand{\Kc}{\mathcal K}
\newcommand{\len}{\operatorname{len}}
\newcommand{\ord}{\operatorname{ord}}
\newcommand{\tl}[1]{\texttt{#1}}
\setlength{\emergencystretch}{3em}

\title{Frames at odd positions and the unit scalars}
\author{}
\date{November 26, 2026}

\begin{document}
\maketitle
\refstepcounter{chap}\label{chap}

\paragraph{Prerequisites.}
Chapter~\ref{di-chap}, Section~\ref{di-sec:chain}, with conditions (a)--(c) of Theorem~\ref{di-thm:true-lifts}, which are Adriazola's hypotheses (N3), (N1) and (N2) \cite[Hyps.~11.42, 11.22 and~11.38]{Adr}, and Chapter~\ref{nk-chap}, which proves (a) for the kernels of the Lean formalization.  We use the data at an inert prime of \cite[\S11.1]{Adr}: the CM field \(E\supset K\) with CM type \(\Phi\), a prime \(p\ge5\) inert in \(K\) and unramified in \(E\), \(W\), \(\sigma\), \(\Kc_0\) and \(\iota\) as in Chapter~\ref{nk-chap}, and \(\ord_p\) with \(\ord_p(p)=1\).  Linear algebra over a discrete valuation ring is all that is needed.

Hypothesis (N2), condition (c) of Theorem~\ref{di-thm:true-lifts}, asks that at the odd positions of an orbit the passage from bounds on kernel entries to inclusions of lattices works as it does at the even positions, with a determinant factor of the opposite sign.  Hypothesis (N1), condition (b), asks that the scalar normalizing each step kernel change only by units when the lattice at one orbit changes.  Why are these not automatic, and why does the Lean formalization \cite{Lean} not state (N2) at all?  The answer to the second question is a choice of bases: at an odd position the Lean uses the dual of the basis at the preceding even position.

\section{Positions and Frobenius}\label{sec:positions}
An \emph{orbit} is a place \(v\) of the maximal real subfield \(E^+\) above \(p\) that splits in \(E\).  Fix \(w\mid v\) as in \cite[\S11.1]{Adr}, put \(Q_v=E_w\), which is unramified of degree \(2\nu\) over \(\Q_p\), let \(W_v\) be its ring of integers, and fix an embedding \(\rho_v\colon Q_v\to\Kc_0\) extending \(\iota|_K\).  The \emph{positions} of \(v\) are \(s\in\Z/2\nu\), and to each corresponds \(\tau_{v,s}\in\Phi\) with
\[
 \iota\tau_{v,s}(x)=\sigma^s\rho_v(x)\ \ (s\ \text{even}),\qquad
 \iota\tau_{v,s}(x)=\sigma^s\rho_v(\bar x)\ \ (s\ \text{odd})
\]
for \(x\in E\).  Even positions see the completion at \(w\), odd positions the completion at \(\bar w\).

\begin{problem}\label{pr:frob}
For every \(s\in\Z/2\nu\) and \(x\in E\), \(\sigma\big(\iota\tau_{v,s}(x)\big)=\iota\tau_{v,s+1}(\bar x)\).  In particular \(\sigma^2\circ\iota\tau_{v,s}=\iota\tau_{v,s+2}\).
\end{problem}

So Frobenius composed with complex conjugation moves each position to the next.  The number of positions is even because \(Q_v\) contains the completion \(\Q_{p^2}\) of \(K\).

\section{Dual frames}\label{sec:frames}
For a \(W_v\)-lattice \(\Lambda\subset Q_v^N\) with basis matrix \(g\in\mathrm{GL}_N(Q_v)\) put \(\len\Lambda=-\ord_p\det g\), and let \(\Lambda^*=\{y\in Q_v^N:\ y^{\mathsf T}x\in W_v\text{ for all }x\in\Lambda\}\).

\begin{problem}\label{pr:dual-lattice}
\(g^{-\mathsf T}\) is a basis matrix of \(\Lambda^*\).  Hence \(\Lambda^{**}=\Lambda\) and \(\len\Lambda^*=-\len\Lambda\).
\end{problem}

Let \(N\subset E^n\) be an \(\Oc_E\)-lattice that is self-dual for \(h_n(x,y)=\sum_ix_i\bar y_i\).  At the orbit, \(N\otimes\Z_p\) has the two components \(N_w\subset Q_v^n\) and \(N_{\bar w}\), and self-duality says that \(N_{\bar w}\) is the image of \(N_w^*\) under complex conjugation.  By the formula for \(\iota\tau_{v,s}\), a basis matrix \(g\) of \(N_w\) therefore gives the frame
\[
 g_s=\sigma^s\rho_v(g)\ \ (s\ \text{even}),\qquad g_s=\sigma^s\rho_v\big(g^{-\mathsf T}\big)\ \ (s\ \text{odd})
\]
of \(N\) at the position \(s\).

\begin{problem}\label{pr:frames}
\(\sigma(g_s)=g_{s+1}^{-\mathsf T}\) and \(\sigma^2(g_s)=g_{s+2}\) for every \(s\in\Z/2\nu\).  Moreover \(\ord_p\det g_s=-\len N_w\) for even \(s\) and \(\ord_p\det g_s=\len N_w\) for odd \(s\).
\end{problem}

The Lean states the same relation globally.  It chooses, for each \(\tau\in\Phi\), a matrix \(P_\tau\in\mathrm{GL}_n(E)\) whose columns form a basis of \(N\) at the place induced by \(\tau\), subject to
\[
 \overline{P_{\tau_{v,s+1}}}^{\,\mathsf T}\,P_{\tau_{v,s}}=1\qquad\text{for all orbits \(v\) and positions \(s\)}.
\]

\begin{problem}\label{pr:dual-basis}
Let \(F\) be a field with an involution \(x\mapsto\bar x\), and \(h(x,y)=x^{\mathsf T}\bar y\) on \(F^n\).  For \(P,Q\in\mathrm{GL}_n(F)\), \(\bar Q^{\mathsf T}P=1\) if and only if \(h(Pe_i,Qe_j)=[i=j]\) for all \(i,j\), that is, the columns of \(Q\) are the \(h\)-dual basis of the columns of \(P\).  For \(F=E\) with complex conjugation, if \(\bar Q^{\mathsf T}P=1\), then \(g_s=\iota\tau_{v,s}(P)\) and \(g_{s+1}=\iota\tau_{v,s+1}(Q)\) satisfy \(\sigma(g_s)=g_{s+1}^{-\mathsf T}\).
\end{problem}

\begin{problem}\label{pr:unitary-dual}
If \(\gamma\in\mathrm{GL}_n(E)\) is unitary, \(\bar\gamma^{\mathsf T}\gamma=1\), then \(\bar\gamma^{-\mathsf T}=\gamma\).  Hence \((P_\tau)\mapsto(\gamma P_\tau)\) preserves the relation displayed above.
\end{problem}

Problem~\ref{pr:unitary-dual} is why the unitary group acts on the frames at all positions by the same matrix, as in the covariance of the lifts \cite[Lemma~11.11]{Adr}: at an odd position the contragredient action \(g\mapsto g^{-\mathsf T}\) of the model is the ordinary action on the dual frame.

\section{The odd positions, (N2)}\label{sec:N2}
Recall from \cite[Def.~11.3, Lemma~11.4]{Adr} that the integral structure \(H_s(\Lambda)\) at a position is the image of the minimal lattice under the frame, twisted by a power of its determinant, and that an edge \((\Lambda'',\Lambda)\) of type \(i\), meaning \(p\Lambda\subset\Lambda''\subset\Lambda\) with \(\Lambda/\Lambda''\) of dimension \(i\), has a basis \(g\) of \(\Lambda\) with \(gd\) a basis of \(\Lambda''\), where \(d=\operatorname{diag}(p,\dots,p,1,\dots,1)\) has \(i\) entries \(p\).

\begin{problem}\label{pr:edge-dual}
In this situation \(\Lambda''^*\) has basis matrix \(g^{-\mathsf T}d^{-1}\), and \(p\Lambda''^*\subset\Lambda^*\subset\Lambda''^*\) with \(\Lambda''^*/\Lambda^*\) of dimension \(i\).
\end{problem}

A \emph{canonical bideterminant} in an \(n\times n\) matrix \(X\) of indeterminates is a product \(B=\prod_c\Delta_{S_c}\), where \(\Delta_S\) is the minor on the rows \(S\) and the first \(\#S\) columns \cite[Def.~11.1]{Adr}.

\begin{problem}\label{pr:bidet}
For \(D=\operatorname{diag}(p^{a_1},\dots,p^{a_n})\), \(B(DX)=p^{\sum_c\sum_{r\in S_c}a_r}B(X)\).  In particular \(B(d^{\pm1}X)=p^{\pm\sum_c\#(S_c\cap\{1,\dots,i\})}B(X)\), and \(0\le\#(S_c\cap\{1,\dots,i\})\le\min(\#S_c,i)\).
\end{problem}

At an even position the frame of the edge lattice is \(\sigma^s\rho_v(gd)\).  The change of frame divides the bideterminants by \(p^{\sum_c\#(S_c\cap\{1,\dots,i_0\})}\), and the normalizing power \(p^{\sum_{a\le i_0}\lambda(a)}\) of the Hecke operator leaves the exponent \(\sum_c\big(\min(h'_c,i_0)-\#(S_c\cap\{1,\dots,i_0\})\big)\ge0\) \cite[Lemma~11.4, Thm.~11.35]{Adr}.  At an odd position the frame of the edge lattice is \(\sigma^s\rho_v(g^{-\mathsf T}d^{-1})\) by Problem~\ref{pr:edge-dual}.  The diagonal matrix enters with the opposite exponent, and by Problem~\ref{pr:bidet} the bideterminants are multiplied by \(p^{\sum_c\#(S_c\cap\{1,\dots,i_0\})}\) instead.  This is the exponent \(\#(S_c\cap\{1,\dots,i_0\})\) in part (a) of (N2) \cite[Hyp.~11.38]{Adr}.  Part (b) concerns the determinant twists.

\begin{problem}\label{pr:twist}
Let \(N_w\) and \(N'_w\) have lengths \(\ell\) and \(\ell'\), with frames \(g_s,g'_s\) as in Problem~\ref{pr:frames}, and let \(w_s\) be integers.  Then
\[
 \ord_p\prod_s\Big(\frac{\det g_s}{\det g'_s}\Big)^{w_s}=\mu\,(\ell'-\ell),\qquad \mu=\sum_{s\ \mathrm{even}}w_s-\sum_{s\ \mathrm{odd}}w_s .
\]
At an odd position \(\big(\det g_s/\det g'_s\big)^{w_s}=\sigma^s\rho_v\big((\det g/\det g')^{-w_s}\big)\).
\end{problem}

With \(w_s=w_{j+1}(v,s)\), the number \(\mu\) is \(\mu_v(j+1)\), and the product cancels the factor \(p^{\mu_v(j+1)(\len L_v-\len L'_v)}\) in the step map \(\kappa_j\) \cite[Def.~11.9]{Adr}.  The last assertion of Problem~\ref{pr:twist} is part (b) of (N2).

\emph{Status}: Adriazola states (N2) as a hypothesis and says that the bookkeeping of the determinant powers at the odd positions remains to be written \cite[Hyp.~11.38, Remark~11.39]{Adr}.  Problems~\ref{pr:edge-dual}--\ref{pr:twist} are that bookkeeping for the frames and the edges.  The Lean formalization has no separate statement.  Its frames at odd positions are the dual frames of Section~\ref{sec:frames}, and the lattices of the chain are described at every position, even or odd, by the integrality of the pairings of the coordinates with one family of products of minors, so both parities are handled by one argument.  We have not checked in detail how this description at the odd positions matches the dual model weights of \cite[\S11.1]{Adr}.

\section{The unit scalars, (N1)}\label{sec:N1}
Each step kernel carries a scalar \(s_j(c,L;c',L')\), built from a step character \(\mu_{j+1}\) of the norm-one ideles, a frame factor and powers of a Hecke character \(\chi\) of \(E\) \cite[Def.~11.23]{Adr}.  (N1) asks that when the lattice changes at one orbit \(v\), the scalar change by \(a_v^{\Delta\len}a_v'^{\Delta\len'}\) with units \(a_v,a'_v\).  Let \(\mathfrak P_w,\mathfrak P_{\bar w}\) be the primes of \(\Oc_E\) at \(w,\bar w\).  The elements \(\tau\in\Phi\) that induce \(w\) or \(\bar w\) are exactly the \(\tau_{v,s}\); every other \(\tau\in\Phi\) induces a place above \(p\) other than \(w,\bar w\) \cite[proof of Lemma~11.25]{Adr}.

\begin{problem}\label{pr:pi-v}
The idele \(\pi_v\) with component \(p\) at \(w\), \(p^{-1}\) at \(\bar w\) and \(1\) elsewhere satisfies \(\pi_v\bar\pi_v=1\), and its ideal is \(\mathfrak P_w\mathfrak P_{\bar w}^{-1}\).
\end{problem}

\begin{problem}\label{pr:chi-count}
Let \(k\ge1\) and \(\beta\in E\) with \(\beta\bar\beta=1\) and \(\beta\Oc_E=(\mathfrak P_{\bar w}\mathfrak P_w^{-1})^k\).  Then \(\ord_p\iota\tau_{v,s}(\beta)\) is \(-k\) for even \(s\) and \(k\) for odd \(s\), and \(\sum_{\tau\in\Phi}\ord_p\iota\tau(\beta)=0\).
\end{problem}

Such \(\beta\) exist for \(k\) twice the class number of \(E\), and for them \(\chi(\mathfrak P_{\bar w}\mathfrak P_w^{-1})^k=\prod_{\tau\in\Phi}\tau(\beta)^{-1}\); so Problem~\ref{pr:chi-count} shows that \(\iota\chi(\mathfrak P_{\bar w}\mathfrak P_w^{-1})\) is a unit \cite[Lemma~11.30]{Adr}.

\begin{problem}\label{pr:mu-val}
Let \(k\ge1\) and \(z\in E\) with \(z\bar z=1\) and \(z\Oc_E=(\mathfrak P_w\mathfrak P_{\bar w}^{-1})^{k}\), let \(\varepsilon_\tau\) \((\tau\in\Phi)\) and \(c\) be integers, and put \(w_s=\varepsilon_{\tau_{v,s}}-c\).  Then \(\sum_{\tau\in\Phi}\varepsilon_\tau\ord_p\iota\tau(z)=k\big(\sum_{s\ \mathrm{even}}w_s-\sum_{s\ \mathrm{odd}}w_s\big)\).
\end{problem}

If \(\mu_{j+1}(z)=\prod_\tau\tau(z)^{\varepsilon_\tau}\) on norm-one elements and \(\mu_{j+1}\) is trivial on the local units at \(v\), then \(\mu_{j+1}(\pi_v^{k}z^{-1})\) is a root of unity, and Problem~\ref{pr:mu-val} gives \(\ord_p\iota\mu_{j+1}(\pi_v)=\mu_v(j+1)\) \cite[Lemma~11.25]{Adr}.  With the effect of changing a lattice on the determinant idele \cite[Lemma~11.26]{Adr}, this proves (N1) for Adriazola's normalization, with
\(a_v=p^{-\mu_v(j+1)}\iota\mu_{j+1}(\pi_v)\,\iota\chi((\pi_v))^{a_j}\) and \(a'_v=p^{\mu_v(j+1)}\iota\mu_{j+1}(\pi_v)^{-1}\iota\chi((\pi_v))^{-a_{j+1}}\) \cite[Prop.~11.27]{Adr}.

\emph{Status}: (N1) for the normalization of \cite[Def.~11.23]{Adr} is proved in \cite[Prop.~11.27]{Adr}.  The Lean formalization writes each step map in local coordinates as a scalar times the pairing of the source coordinates with the kernel series in adapted frames, and proves that the scalar is a \(p\)-adic unit.  We have not followed that proof, and we have not compared the Lean's scalar with \(s_j\) term by term.

\begin{leanremark}\sloppy
The dual frames are \tl{GlobalTheory.KerStd.StdBases}, whose second clause is the relation \(\overline{P_{\tau_{v,s+1}}}^{\mathsf T}P_{\tau_{v,s}}=1\).  Problem~\ref{pr:frob} is \tl{exists\_frob\_step}, and Problems~\ref{pr:frames} and~\ref{pr:dual-basis} are \tl{frob\_moves\_coords}.  The description of the lattices by pairings, at all positions, is in the fields \tl{lat\_S} and \tl{levelFree\_lattice\_condition} of the structure \tl{KerEnt}; its field \tl{kform} writes the step map as \tl{un} times the pairing with the kernel.  The unit property of \tl{un} is \tl{KerEnt.UnUnit}, proved by \tl{GlobalTheory.T.kernelPairScalars\_isUnit\_trueColumns}.  The theorem \tl{GlobalTheory.PdvClose.at\_of\_true\_columns} combines it with the kernel law of Chapter~\ref{nk-chap} \cite{Lean}.
\end{leanremark}

\section*{Hints}
\noindent\emph{Section~\ref{sec:positions}.}
\ref{pr:frob}. Treat even and odd \(s\) separately; \(\sigma^{2\nu}\) fixes \(\rho_v(Q_v)\).

\noindent\emph{Section~\ref{sec:frames}.}
\ref{pr:dual-lattice}. \(y\in\Lambda^*\) if and only if \(g^{\mathsf T}y\in W_v^N\).
\ref{pr:frames}. The definition of \(g_s\), and \(\sigma^{2\nu}\rho_v=\rho_v\) as in Problem~\ref{pr:frob}.
\ref{pr:dual-basis}. Transpose; then Problem~\ref{pr:frob}.

\noindent\emph{Section~\ref{sec:N2}.}
\ref{pr:edge-dual}. Problem~\ref{pr:dual-lattice} for \(\Lambda''\).
\ref{pr:bidet}. Multiplying a row by \(p^{a}\) multiplies every minor containing it by \(p^{a}\).
\ref{pr:twist}. Problem~\ref{pr:frames}.

\noindent\emph{Section~\ref{sec:N1}.}
\ref{pr:chi-count}. The formula for \(\iota\tau_{v,s}\), and \(\ord_w(\bar\beta)=\ord_{\bar w}(\beta)\); there are \(\nu\) positions of each parity.
\ref{pr:mu-val}. As in Problem~\ref{pr:chi-count}.

\begin{thebibliography}{Lean}
\bibitem[Adr]{Adr}
J.~Adriazola, \emph{The Alp\"oge proof of the Riemann hypothesis}, manuscript, revised 24 November 2026, \url{https://github.com/jimjamzola/RH-Alpoge-Reconstruction}.
\bibitem[Lean]{Lean} The Lean community, \emph{riemann-hypothesis}, a Lean formalization of Alp\"oge's argument, \url{https://github.com/leanprover-community/riemann-hypothesis}.
\end{thebibliography}
\end{document}
